\section{Analysis plan and estimator notes}
\label{app:b-analysis}

\emph{Status.} This appendix records the analysis conventions the first study
uses and points at the historical full plan for the deferred E3 design.  No
analysis here has been executed on agent data.

\subsection{Units, clustering, and folds}

Roots instantiated from one task template are dependent, so template identity
induces within-cluster dependence that the design does not remove.  Task
templates are grouped before splitting, so two instances of one template cannot
straddle a fold; the clustering level fixes the resampling unit of the cluster
bootstrap, and the effective cluster count follows from that level rather than
from the root count.  Every distribution-free statement for $n$ independent
units is a statement about independent \emph{templates}, and reading $n$ as the
root count is anticonservative.  Splits are training, selection, calibration and
an untouched final evaluation set, frozen after the pilot.

\subsection{Minimum detectable effects}

For a within-checkpoint contrast of a binary outcome with $R$ repetitions per
arm and between-checkpoint standard deviation $\sigma_b$, the minimum detectable
effect at $80\%$ power and two-sided $5\%$ level is
$(z_{0.975}+z_{0.80})\sqrt{(\sigma_b^2+0.5/R)/n}$.  This is a generic planning
baseline; the operative final-test precision is computed from pilot variance
with conservative uncertainty and recorded before outcomes.

\subsection{Risk-ranking regret: the registered estimand}
\label{app:e3-estimator}

For a capacity $q\in(0,1)$ and a score $s$, let $S_q(s)$ denote the upper
$q$-fraction of the eligible population under $s$.  With $a_v(w)$ the
verification advantage on the common-availability domain, the budget-matched
risk-ranking regret is
\begin{equation}
 \mathcal R_{\mathrm{rank}}(q)
 =\E\big[a_v(W)\mid W\in S_q(a_v)\big]-\E\big[a_v(W)\mid W\in S_q(r)\big].
 \label{eq:b-rank-regret}
\end{equation}
The registered grid is $q\in\{0.05,0.10,0.25,0.50\}$; the effect score is fitted
on template-grouped folds and evaluated on independent repetitions.  The
forbidden plug-in procedure --- using the same noisy per-checkpoint estimate to
select and evaluate --- is a selection-on-noise estimator and is not used; the
estimator audit reproduces the large positive bias it induces under a
simulated zero-regret null.  The full historical specification, including the
retired F1 equivalence design and the F2-FEHPO power report, is preserved in
the repository; this appendix does not re-register it.  The first study's
primary null is the frozen-learned-policy comparison of
\cref{sec:b-rejection}, and $\mathcal R_{\mathrm{rank}}$ is reported as a
population-level descriptive target.

\subsection{Descriptive action-value baseline}
\label{app:label-use}

The plain action-conditioned baseline (f) of \cref{sec:policy} regresses the
branch outcome on the controller's observation $(Z,M)$ and the regime.  It is a
coherent deployable learner and is reported as such; a regression on the full
state $X$ would estimate the action value itself but is an oracle diagnostic,
because $X$ carries the latent stratum and post-intervention state.  Neither is
the masked policy-score learner of \cref{eq:rl-objective}.

\subsection{Proposed extensions}
\label{app:extensions}

Four extensions are specified outside this manuscript in
\texttt{docs/paper\_track/} and are named so the boundary between them and the
first study is explicit: a prospective verifier-information study separating
the information a verifier supplies from the response it triggers; a
matching-diagnostic study for coarse-matching distortion; a calibration-size
sensitivity study; and the recovery extension of \cref{sec:b-bridge}.  Each
carries its own candidate registry and none is part of the first study's
confirmatory claims.
